\chapter{Phénomènes énergétiques : respiration cellulaire et fermentation}

\noindent\textit{Pour fonctionner, se contracter ou simplement rester en vie, les cellules
ont besoin d'énergie. Cette énergie provient de la dégradation du glucose, selon deux voies
possibles : la respiration cellulaire, qui utilise le dioxygène, et la fermentation, qui
s'en passe.}
\vspace{8pt}

\connecteBanner

\section{La respiration cellulaire}

\begin{definition}[Respiration cellulaire]
La \textbf{respiration cellulaire} est la réaction par laquelle une cellule dégrade
complètement le \textbf{glucose} en présence de \textbf{dioxygène}, pour libérer de
l'\textbf{énergie} utilisable :
\[
\mathrm{C_6H_{12}O_6 + 6\,O_2 \longrightarrow 6\,CO_2 + 6\,H_2O + \text{énergie}}.
\]
\end{definition}

\begin{propriete}[Origine des réactifs]
Le \textbf{glucose} provient de la digestion des aliments, transporté par le sang depuis
l'intestin. Le \textbf{dioxygène} provient de l'air, capté au niveau des alvéoles
pulmonaires (hématose, voir chapitre précédent) et transporté par le sang jusqu'aux
cellules.
\end{propriete}

\begin{illus}
\begin{tikzpicture}[scale=0.8]
\filldraw[onbo!20,draw=onbo,line width=1.1pt] (0,0) rectangle (2.6,0.9);
\node[font=\small] at (1.3,0.45) {glucose $+$ $\mathrm{O_2}$};
\draw[->,onbgreen,line width=1.3pt] (2.6,0.45) -- (4.3,0.45);
\node[above,onbgreen,font=\scriptsize] at (3.45,0.65) {respiration};
\filldraw[onbgreen!20,draw=onbgreen,line width=1.1pt] (4.8,0) rectangle (8.7,0.9);
\node[font=\small] at (6.75,0.45) {$\mathrm{CO_2}$ $+$ $\mathrm{H_2O}$ $+$ énergie};
\end{tikzpicture}
\fig{Figure — Bilan simplifié de la respiration cellulaire.}
\end{illus}

\section{La fermentation}

\begin{definition}[Fermentation]
La \textbf{fermentation} est une dégradation \textbf{incomplète} du glucose, qui ne
nécessite \textbf{pas de dioxygène}. Elle libère de l'énergie, mais en quantité
\textbf{beaucoup plus faible} que la respiration, pour une même quantité de glucose
dégradée.
\end{definition}

\begin{propriete}[Fermentation lactique — muscle en effort intense]
Lorsque l'apport de dioxygène devient insuffisant (effort musculaire intense), les
cellules musculaires basculent vers la fermentation lactique :
\[
\mathrm{C_6H_{12}O_6 \longrightarrow 2\,C_3H_6O_3 + \text{énergie}}
\]
(le produit formé, $\mathrm{C_3H_6O_3}$, est l'\textbf{acide lactique}).
\end{propriete}

\begin{remarque}
L'accumulation d'acide lactique dans le muscle est responsable de la sensation de
brûlure et de fatigue musculaire ressentie lors d'un effort intense et prolongé.
\end{remarque}

\begin{propriete}[Fermentation alcoolique — levures]
Certains micro-organismes, comme les levures, réalisent une fermentation
\textbf{alcoolique} :
\[
\mathrm{C_6H_{12}O_6 \longrightarrow 2\,C_2H_6O + 2\,CO_2 + \text{énergie}}
\]
(le produit $\mathrm{C_2H_6O}$ est l'\textbf{éthanol}). Cette réaction est utilisée dans
la fabrication du pain (le $\mathrm{CO_2}$ fait lever la pâte) et des boissons
fermentées.
\end{propriete}

\section{Comparer respiration et fermentation}

\begin{propriete}[Tableau comparatif]
\begin{center}
\small
\begin{tabular}{lcc}
\toprule
& \textbf{Respiration} & \textbf{Fermentation} \\
\midrule
Dioxygène nécessaire & Oui & Non \\
Dégradation du glucose & Complète & Incomplète \\
Énergie libérée & Élevée & Faible \\
Déchets & $\mathrm{CO_2}$, $\mathrm{H_2O}$ & acide lactique, ou éthanol $+$ $\mathrm{CO_2}$ \\
\bottomrule
\end{tabular}
\end{center}
\end{propriete}

\begin{exemple}
Un sprinteur qui court à allure maximale ne peut pas apporter assez de dioxygène à ses
muscles : une partie de l'énergie provient alors de la fermentation lactique, ce qui
explique les courbatures ressenties le lendemain.
\end{exemple}

% ═══════════════════════════════════════════════════════════════
\section{L'essentiel à retenir}

\begin{synthese}
\textbf{Respiration cellulaire.} $\mathrm{C_6H_{12}O_6+6\,O_2\to6\,CO_2+6\,H_2O+}$
énergie ; dégradation complète, énergie élevée.\\[4pt]
\textbf{Fermentation lactique.} $\mathrm{C_6H_{12}O_6\to2\,C_3H_6O_3+}$ énergie (muscle,
sans $\mathrm{O_2}$).\\[4pt]
\textbf{Fermentation alcoolique.} $\mathrm{C_6H_{12}O_6\to2\,C_2H_6O+2\,CO_2+}$ énergie
(levures).\\[4pt]
\textbf{Comparaison.} La fermentation libère beaucoup moins d'énergie que la respiration,
pour une même quantité de glucose.\\[4pt]
\textbf{Piège fréquent.} La fermentation n'est pas « meilleure » ou « plus rapide » : elle
est simplement la seule option quand le dioxygène manque, avec un rendement énergétique
bien plus faible.
\end{synthese}

% ═══════════════════════════════════════════════════════════════
\section{Méthodes \& savoir-faire}

\methode{Écrire et vérifier l'équation de la respiration cellulaire}
Partir de glucose $+$ dioxygène, produire dioxyde de carbone $+$ eau $+$ énergie, puis
compter les atomes de chaque côté.
\begin{exemple}
$\mathrm{C_6H_{12}O_6+6\,O_2\to6\,CO_2+6\,H_2O}$ : vérifier C $6=6$, H $12=12$, O
$18=6\times2+6=18$.
\end{exemple}

\methode{Distinguer respiration et fermentation}
Se demander si le dioxygène est disponible en quantité suffisante : oui $\to$
respiration ; non $\to$ fermentation.
\begin{exemple}
Une pâte à pain qui lève grâce aux levures, sans air, réalise une fermentation
alcoolique.
\end{exemple}

\methode{Interpréter une situation d'effort musculaire}
Identifier si l'apport en dioxygène est suffisant (respiration) ou dépassé par la demande
(fermentation lactique, acide lactique, fatigue musculaire).
\begin{exemple}
Une crampe après un sprint s'explique par une accumulation d'acide lactique issue de la
fermentation lactique.
\end{exemple}

% ═══════════════════════════════════════════════════════════════
\section{Exercices d'application}
\begin{multicols}{2}

\rubrique{Respiration cellulaire}
\exo{}\diff{1} Écrire l'équation-bilan simplifiée de la respiration cellulaire.

\exo{}\diff{2} D'où proviennent le glucose et le dioxygène utilisés par les cellules ?

\exo{}\diff{2} Quels sont les produits de la respiration cellulaire ?

\rubrique{Fermentation}
\exo{}\diff{1} Qu'est-ce que la fermentation ?

\exo{}\diff{2} Quel produit forme la fermentation lactique ?

\exo{}\diff{2} Quels produits forme la fermentation alcoolique ?

\exo{}\diff{2} Dans quelle situation le muscle utilise-t-il la fermentation lactique
plutôt que la respiration ?

\rubrique{Comparaison}
\exo{}\diff{2} Laquelle des deux voies, respiration ou fermentation, nécessite du
dioxygène ?

\exo{}\diff{2} Laquelle des deux voies libère le plus d'énergie pour une même quantité de
glucose ?

\exo{}\diff{3} Pourquoi ressent-on des courbatures après un effort physique intense et
inhabituel ?

% ═══════════════════════════════════════════════════════════════
\end{multicols}

\section{Exercices d'entraînement}
\begin{multicols}{2}

\rubrique{Respiration cellulaire}
\exo{}\diff{2} Vérifier que l'équation-bilan de la respiration cellulaire, donnée plus
haut, est bien équilibrée (compter les atomes de C, H et O de chaque côté).

\exo{}\diff{3} Expliquer le lien entre la respiration pulmonaire (chapitre précédent) et
la respiration cellulaire.

\rubrique{Fermentation}
\exo{}\diff{2} Vérifier que l'équation de la fermentation lactique
$\mathrm{C_6H_{12}O_6\to2\,C_3H_6O_3}$ est équilibrée.

\exo{}\diff{2} Vérifier que l'équation de la fermentation alcoolique
$\mathrm{C_6H_{12}O_6\to2\,C_2H_6O+2\,CO_2}$ est équilibrée.

\exo{}\diff{3} Pourquoi la fabrication du pain utilise-t-elle la fermentation alcoolique
des levures ?

\rubrique{Comparaison}
\exo{}\diff{3} Un athlète de fond (longue distance) et un sprinteur (course très courte
et intense) sollicitent-ils les deux voies énergétiques dans les mêmes proportions ?
Justifier.

\exo{}\diff{3} Pourquoi la respiration est-elle qualifiée de dégradation « complète » du
glucose, contrairement à la fermentation ?

\exo{}\diff{2} Citer un point commun entre la respiration cellulaire et la fermentation
(en dehors du glucose de départ).

\exo{}\diff{3} Expliquer pourquoi un muscle ne peut pas fonctionner indéfiniment par
fermentation lactique seule.

% ═══════════════════════════════════════════════════════════════
\end{multicols}

\section{Sujets type BEPC}

\sujetbac[La respiration cellulaire]
\begin{enumerate}[label=\arabic*)]
\item Écris l'équation-bilan simplifiée de la respiration cellulaire.
\item D'où proviennent le glucose et le dioxygène nécessaires à cette réaction ?
\item Quels sont les produits formés ?
\item Pourquoi cette réaction est-elle indispensable au fonctionnement des cellules ?
\end{enumerate}

\sujetbac[Fermentation et effort musculaire]
\begin{enumerate}[label=\arabic*)]
\item Qu'est-ce que la fermentation lactique ?
\item Dans quelle situation un muscle a-t-il recours à cette voie ?
\item Quel produit s'accumule dans le muscle, et quelle sensation cela provoque-t-il ?
\item Compare, en une phrase, l'énergie libérée par la fermentation à celle libérée par
la respiration.
\end{enumerate}

\sujetbac[La fermentation alcoolique et ses usages]
\begin{enumerate}[label=\arabic*)]
\item Écris l'équation-bilan de la fermentation alcoolique.
\item Quels organismes réalisent cette réaction ?
\item Pourquoi cette réaction est-elle utilisée dans la fabrication du pain ?
\item Cette réaction nécessite-t-elle du dioxygène ? Justifie.
\end{enumerate}

% ═══════════════════════════════════════════════════════════════
\section{Corrigés}
\begin{multicols}{2}

\corrige{de l'exercice 1}
$\mathrm{C_6H_{12}O_6+6\,O_2\to6\,CO_2+6\,H_2O+}$ énergie.

\corrige{de l'exercice 2}
Le glucose provient de la digestion des aliments ; le dioxygène provient de l'air, capté
au niveau des alvéoles pulmonaires.

\corrige{de l'exercice 3}
Le dioxyde de carbone, l'eau, et de l'énergie.

\corrige{de l'exercice 4}
Une dégradation incomplète du glucose, sans dioxygène, libérant peu d'énergie.

\corrige{de l'exercice 5}
L'acide lactique.

\corrige{de l'exercice 6}
De l'éthanol et du dioxyde de carbone.

\corrige{de l'exercice 7}
Lors d'un effort intense, quand l'apport de dioxygène devient insuffisant pour répondre
aux besoins du muscle.

\corrige{de l'exercice 8}
La respiration.

\corrige{de l'exercice 9}
La respiration.

\corrige{de l'exercice 10}
Car les muscles ont utilisé la fermentation lactique, qui a fait s'accumuler de l'acide
lactique responsable de la sensation douloureuse.

\corrige{de l'exercice 11}
Gauche : C $=6$, H $=12$, O $=6+12=18$. Droite : C $=6$, H $=12$, O $=12+6=18$. Équation
équilibrée.

\corrige{de l'exercice 12}
Gauche : C $=6$, H $=12$, O $=6$. Droite : C $=2\times3=6$, H $=2\times6=12$, O
$=2\times3=6$. Équation équilibrée.

\corrige{de l'exercice 13}
Gauche : C $=6$, H $=12$, O $=6$. Droite : C $=2\times2+2=6$, H $=2\times6=12$, O
$=2\times1+2\times2=6$. Équation équilibrée.

\corrige{de l'exercice 14}
Car le dioxyde de carbone produit par les levures forme des bulles qui font lever la
pâte.

\corrige{du sujet 1}
1) $\mathrm{C_6H_{12}O_6+6\,O_2\to6\,CO_2+6\,H_2O+}$ énergie.\\
2) Le glucose vient de la digestion des aliments ; le dioxygène vient de l'air, capté aux
alvéoles pulmonaires.\\
3) Le dioxyde de carbone, l'eau et de l'énergie.\\
4) Car c'est cette énergie qui permet aux cellules de fonctionner (contraction
musculaire, etc.).

\corrige{du sujet 2}
1) Une dégradation incomplète du glucose en acide lactique, sans dioxygène.\\
2) Lorsque l'effort musculaire est trop intense pour que l'apport en dioxygène suffise.\\
3) L'acide lactique s'accumule, provoquant une sensation de brûlure et de fatigue
musculaire.\\
4) La fermentation libère beaucoup moins d'énergie que la respiration, pour une même
quantité de glucose.

\corrige{du sujet 3}
1) $\mathrm{C_6H_{12}O_6\to2\,C_2H_6O+2\,CO_2+}$ énergie.\\
2) Les levures (micro-organismes).\\
3) Car le dioxyde de carbone produit fait lever la pâte.\\
4) Non, la fermentation ne nécessite pas de dioxygène : c'est justement ce qui la
distingue de la respiration.

\end{multicols}
\exercicesBanner
